\scanpage{64}
\mspage{55}
% archive sheet number: 64
$1^\circ)$ On veut qu'elle soit croissante pour $\leq$, et commute aux $\mathrm{Sup}$ quelconques. Ça implique qu'elle sera déterminée par sa restriction aux fig él.
\[
\unsure{f_M|f_\mu}:\unsure{M|\mathcal{M}}\longrightarrow\unsure{\mathscr{F}'|\mathcal{F}'|F'},
\]
qui sera une application croissante pour $\leq$. Si les figures de $\mathscr{F}$ sont de type fini, ou celles de $\mathscr{F}'$ \unsure{inconditionnelles}, n'importe quelle appl. \inserted{$f_M$} croissante pour $\leq$ proviendra d'un $f$.

$2^\circ)$ On veut qu'elle commute aux $\mathrm{Inf}$ quelconques, à familles \inserted{d'indices} non vides, et pour commencer
\[
f(F\cap G)=f(F)\cap f(G).\tag*{(*)}
\]
On est ramené aussitôt au cas où $F,G$ sont des multistrates, et la formule devient
\[
f_M(X)\cap f_M(Y)=\mathop{\mathrm{Sup}}_{Z\in\widetilde{X}\cap\widetilde{Y}}f_M(Z)\tag*{$(*_M)$}
\]
(qui rappelle la formule familière pour une figure \unsure{ensembliste}
\[
\unsure{X_i|x_i}\cap X_j=\bigcup_{k\leq i,j}X_{\unsure{n|k}}
\]
).

Cette formule (*) assure la commutation aux $\mathrm{Inf}$ finis. Pour les $\mathrm{Inf}$ infinis, on est ramené au cas des familles \deleted{\unsure{crois}} filtrantes décroissantes. Il n'est
